\chapter{Hors et dans les structures}
\label{chp.struct.mod}

\minitoc

\lettrine{C}{e} chapitre contient les éléments de base de l'étude des
structures. Là où le chapitre précédent concentrait l'étude sur la
satisfaction et la conservation d'énoncés du premier ordre, ce chapitre sera
dédié à des notions liées directement aux structures d'étude.

Nous abordons d'abord les outils incontournables pour construire des
structures~: l'amalgamation et les chaînes. Nous introduisons ensuite les
parties définissables, et explorons divers résultats qui leur sont liés.

\section{Constructions de structure}

Avant d'étudier des méthodes pour construire des structures, revenons sur une
notion explorée dans la \cref{subsec.enrichissement}~: l'enrichissement de
structure.

\subsection{Enrichissement et diagrammes}

Comme la notion d'enrichissement permet de manipuler plus finement des
structures, on souhaite définir des façons efficaces d'enrichir des structures.
Pour cela, on introduit la notion d'enrichissement par constantes.

\begin{definition}[Enrichissement par constantes]
  Soit $\Sign$ une signature, $\varMM$ une $\Sign$-structure et
  $\varPt\subseteq\carMM$. On appelle enrichissement par constantes de $\Sign$
  sur $\varPt$ la signature
  \[\SignP{\varPt}\defeq \Sign\cup\compr{\Cc_a^0}{a \in \varPt}\]
\end{definition}

On considère donc une nouvelle signature, dans laquelle on a ajouté une
constante pour chaque élément de $\varPt$. Cette signature a un sens naturel
dans la structure $\varMM$ à partir de laquelle on a construit l'enrichissement,
ce qui donne un enrichissement de $\varMM$.

\begin{definition}[Enrichissement par constantes d'une structure]
  Soit $\Sign$ une signature, $\varMM$ une $\Sign$-structure et
  $\varPt\subseteq\carMM$. On appelle enrichissement par constante de $\varMM$
  sur $\varPt$ la $\SignP{\varPt}$-structure $\varMM_{\varPt}$ obtenue en
  associant $\Cc_a$ à $a\in\carMM$.
\end{definition}

Revenons maintenant à un énoncé du chapitre précédent~: le
\cref{lem.equivelt.th}. D'après ce lemme, si on a une structure $\varMM$ et
qu'on veut construire $\varMN$ tel que $\varMM\eltEq\varMN$, il nous suffit de
construire un modèle de $\Th(\varMM)$. Maintenant, grâce à la notion
d'enrichissement par constantes, on peut effectuer cela en rajoutant des
interprétations aux éléments de $\varPt$, ce qui nous permet alors de construire
une théorie caractérisant les extensions élémentaires de $\varMM$~: on appelle
cette théorie le diagramme élémentaire de $\varMM$.

\begin{definition}[Diagramme, diagramme élémentaire]
  Soit $\Sign$ une signature et $\varMM$ une $\Sign$-structure. On appelle
  diagramme de $\varMM$ la théorie
  \[\Diag(\varMM) \defeq
  \compr{\varFF\in\Propqf(\SignP{\carMM})}{\varMM_{\carMM}\satisf\varFF}\]
  et le diagramme élémentaire de $\varMM$ la théorie
  \[\DiagEl(\varMM)\defeq
  \compr{\varFF\in\Prop(\SignP{\carMM})}{\varMM_{\carMM}\satisf\varFF}\]
\end{definition}

\begin{remark}
  Le diagramme élémentaire de $\varMM$ peut directement se définir comme
  étant $\Th(\varMM_{\carMM})$.
\end{remark}

\begin{proposition}
  En reprenant les notations précédentes, on a une correspondance entre
  les modèles de $\Diag(\varMM)$ et les $\Sign$-structures $\varMN$ munies d'un
  plongement $\makeFunName{\varMoM}{\varMM}{\varMN}$. On a aussi une
  correspondance entre les modèles de $\DiagEl(\varMM)$ est le
  $\Sign$-structures $\varMN$ munies d'un plongement élémentaire
  $\makeFunName{\varMoM}{\varMM}{\varMN}$.
\end{proposition}

\begin{proof}
  En effet, si on a $\varMN$ une $\Sign$-structure et $\varMoM$ un plongement
  de $\varMM$ dans $\varMN$, alors on peut construire $\varMN_{\carMM}$ une
  $\SignP{\carMM}$-structure en associant à chaque $\Cc_m$ l'élément
  $f(m)$. On prouve que $\carMN_{\carMM}$ satisfait $\Diag(\varMM)$ en vérifiant
  par induction que pour chaque proposition $\varFF$ sans quantificateur vraie
  dans $\varMM_{\carMM}$, $\carMN_{\carMM}\satisf\varFF$. Dans le cas des
  symboles de relation, cela est vrai par hypothèse sur le fait que
  $\varMoM$ est un plongement, et les autres cas passent automatiquement à
  l'induction. Réciproquement, si $\varMN$ est un modèle de $\Diag(\varMM)$,
  alors on peut construire le plongement
  \[\makeFun{\varMoM}{\varMM}{\varMN}{m}{\toMP{\Cc_m}{\varMN}}\]
  Pour vérifier que ce morphisme est bien un plongement, on procède par
  contraposée~: si $\varMM,\vec m\nsatisf\symbR(\vec x)$, alors
  $\varMM,\vec m\satisf\lnot\symbR(\vec x)$ donc
  $\varMN,f(\vec m)\satisf\lnot\symbR(\vec x)$, \latinexpr{i.e.}
  $\varMN,f(\vec m)\nsatisf\symbR(\vec x)$.

  On raisonne de la même façon dans le cas du diagramme élémentaire, mais
  la situation est plus directe.
\end{proof}

\begin{remark}
  On traite souvent de façon interchangeable deux structures $\varMM,\varMN$
  munies d'un plongement $\makeFunName{\varMoM}{\varMM}{\varMN}$ et une
  extension $\varMM\subseteq\varMN$. De même, on identifie généralement
  une extension élémentaire $\varMM\subModElt\varMN$ et un plongement
  élémentaire $\makeFunName{\varMoM}{\varMM}{\varMN}$. La proposition
  précédente peut donc se lire comme le fait que les modèles de $\Diag(\varMM)$
  correspondent aux extensions de $\varMM$, et les modèles de $\DiagEl(\varMM)$
  correspondent aux extensions élémentaires de $\varMM$.
\end{remark}

Cette méthode ressemble fortement à la preuve du théorème de Löwenheim-Skolem
ascendant, dans laquelle on a introduit une famille de taille $\varCk$ de
constantes, pour construire une extension élémentaire suffisamment grande.
La différence est que dans cette preuve, on considère d'un côté une structure
$\varMM$, et de l'autre un ensemble $\varCk$ (un cardinal, pour être plus
précis), et on cherche à construire une structure $\varMN$ conservant à la fois
la structure $\varMM$ et la taille $\varCk$.

Cette idée, et en particulier l'utilisation de la compacité pour construire
un tel modèle, va être réutilisée ici pour construire de nouvelles structures.

\subsection{Constructions par amalgamation et chaîne}

Avec la méthode des diagrammes, on peut donc transformer un problème de
construction d'extension (élémentaire) en un problème de satisfaction de
théorie. La satisfaction de théorie, elle, bénéficie d'un outil central~: le
théorème de compacité. On peut alors combiner les deux pour construire des
extensions (élémentaires) d'une structure tout en ajoutant d'autres conditions.
Cela mène au théorème de plongement élémentaire conjoint et d'amalgamation
élémentaire.

\begin{theorem}
  Les trois propositions suivantes sont vraies~:
  \begin{enumerate}[label=(\roman*)]
  \item Soient $\varMM_1,\varMM_2$ deux $\Sign$-structures telles que
    $\varMM_1\eltEq\varMM_2$, alors il existe une structure $\varMN$ avec
    deux plongements élémentaires
    \[\makeFunName{\varMoM_1}{\varMM_1}{\varMN}\qquad
    \makeFunName{\varMoM_2}{\varMM_2}{\varMN}\]
  \item Soit $(\varMM_i)_{i \in \indSet}$ une famille de $\Sign$-structures toutes
    élémentairement équivalentes deux à deux, alors il existe une structure
    $\varMN$ et une famille de plongements élémentaires
    \[\makeFunName{\varMoM_i}{\varMM_i}{\varMN}\]
  \item Soient $\varMM_0,\varMM_1,\varMM_2$ trois $\Sign$-structures et deux
    plongements élémentaires
    \[\makeFunName{\varMoM_1}{\varMM_0}{\varMM_1}\qquad
    \makeFunName{\varMoM_2}{\varMM_0}{\varMM_2}\]
    alors il existe une $\Sign$-structure $\varMN$ et deux plongements
    élémentaires
    \[\makeFunName{\varMoMB_1}{\varMM_1}{\varMN}\qquad
    \makeFunName{\varMoMB_2}{\varMM_2}{\varMN}\]
    tels que $\varMoMB_1\circ\varMoM_1 = \varMoMB_2\circ\varMoM_2$,
    c'est-à-dire tel que le diagramme suivant commute~:
    \begin{center}
      \begin{tikzcd}
        \varMM_0 \ar[r,"\varMoM_1"]\ar[d,"\varMoM_2"] &
        \varMM_1\ar[d,"\varMoMB_1"]\\
        \varMM_2\ar[r,"\varMoMB_2"] & \varMN
      \end{tikzcd}
    \end{center}
  \end{enumerate}
\end{theorem}

\begin{proof}
  On prouve chaque proposition~:
  \begin{enumerate}[label=(\roman*)]
  \item Pour simplifier la lecture, on considère que
    $\carM{\varMM_1}\cap\carM{\varMM_2}=\vide$ (quitte à dupliquer une des deux
    structures). On définit alors la signature
    \[\Sign' \defeq \SignP{\carM{\varMM_1}}\cup\SignP{\carM{\varMM_2}}\]
    D'après ce qui a été dit plus tôt, l'existence de la structure $\varMN$
    recherchée revient exactement à la cohérence de la théorie
    \[\varT \defeq \DiagEl(\varMM_1)\cup\DiagEl(\varMM_2)\]
    et, par compacité, on peut se limiter à un ensemble fini quelconque de
    propositions. Quitte à prendre une conjonction, il nous suffit de vérifier
    que pour toute $\Sign$-formule $\varFF(\vec x)$ et donnée de paramètres
    $\vec b$ dans $\varMM_2$, il existe un modèle de
    $\DiagEl(\varMM_1)\cup\{\varFF(\vec{\Cc_b})\}$. Mais alors, comme
    $\varMM_1\eltEq\varMM_2$, on sait que
    $\varMM_1\satisf\exists \vec b, \varFF(\vec b)$, donc on peut
    interpréter $\varMM_1$ en une $\Sign'$-structure par $\Cc_b \mapsto b$
    (et les autres constantes $\Cc$ sont interprétées de façon quelconque)
    pour obtenir un modèle de $\DiagEl(\varMM_1)\cup\{\varFF(\vec{\Cc_b})\}$.
    Ainsi, par compacité, on trouve un modèle $\varMN$ tel que décrit
    précédemment.
  \item On procède de la même façon, en prenant
    \[\Sign'\defeq \bigcup_{i \in \indSet}\SignP{\carM{\varMM_i}}\qquad
    \varT\defeq \bigcup_{i \in \indSet}\DiagEl(\varMM_i)\]
    Par compacité, on peut se limiter à des fragments finis, mais alors il
    n'y a qu'un nombre fini de $i \in \indSet$ apparaissant dans les
    formules. On peut donc raisonner par induction sur le nombre de
    $i$ apparaissant~: le cas où les ensembles sont vides est automatique,
    et l'hérédité revient au cas précédemment traité.
  \item D'après nos hypothèses, on peut considérer
    \[\varMM_1' \defeq \varMM_{1,\carM{\varMM_0}} \qquad
    \varMM_2' \defeq \varMM_{2,\carM{\varMM_0}}\]
    Comme $\varMoM_1$ et $\varMoM_2$ sont des plongements élémentaires, on peut
    vérifier que $\varMM_1'\eltEq\varMM_2'$, donc on peut appliquer la première
    proposition démontrée pour trouver $\varMN$ et $\varMoMB_1,\varMoMB_2$,
    qui sont respectivement une $\SignP{\carM{\varMM_0}}$-structure et des
    plongements élémentaires de $\SignP{\carM{\varMM_0}}$-structures. On
    en déduit qu'en particulier, pour tout $m \in \carM{\varMM_0}$~:
    \begin{align*}
      \varMoMB_1(\varMoM_1(m))
      &= \varMoMB_1\left(\toMP{\Cc_{m}}{\varMM_1'}\right)\\
      &= \toMP{\Cc_{m}}{\varMN}\\
      &= \varMoMB_2\left(\toMP{\Cc_{m}}{\varMM_2'}\right)\\
      \varMoMB_1(\varMoM_1(m)) &= \varMoMB_2(\varMoM_2(m))
    \end{align*}
    donc $\varMN$ se réduit en une $\Sign$-structure vérifiant ce qui est
    attendu.\qedhere
  \end{enumerate}
\end{proof}

\begin{remark}
  L'étudiant catégoricien reconnaîtra probablement une situation proche
  respectivement des coproduits et des sommes amalgamées (ou \og poussé en
  avant\fg). Cependant, il faut prendre garde au fait qu'ici, aucune
  universalité n'est établie. Il existe bien une façon de compléter le
  diagramme en un diagramme commutatif, mais cela ne dit rien vis à vis
  d'autres complétions possibles du diagramme commutatif.
\end{remark}

On donne une autre application possible, qui est celle d'étendre des plongements
élémentaires partiels (chose que l'on va donc définir).

\begin{definition}[Plongement élémentaire partiel]
  Soit $\Sign$ une signature et $\varMM,\varMN$ deux $\Sign$-structure. Un
  plongement élémentaire partiel est la donnée d'un sous-ensemble
  $\varPt\subseteq\varMM$ et d'une fonction
  $\makeFunName{\varMoM}{\varPt}{\varMN}$
  telle que pour toute formule $\varFF\in\Formula(\Sign)$ et donnée de
  paramètres $\vec a$ dans $\varPt$, on a l'équivalence
  \[\varMM,\vec a\satisf\varFF \iff \varMN,\varMoM(\vec a)\satisf\varFF\]
\end{definition}

\begin{remark}
  Lorsqu'il existe un plongement élémentaire partiel, $\varMM$ et $\varMN$ sont
  en particulier élémentairement équivalente. En effet, quitte à ne considérer
  que les propositions et à ignorer les paramètres dans $\varPt$, on établit
  aussi que $\varMM\satisf\varFF\iff\varMN\satisf\varFF$.
\end{remark}

\begin{remark}
  Si $(\varPt,\varMoM)$ est un plongement élémentaire partiel entre
  $\varMM$ et $\varMN$, alors $\varMoM$ induit un isomorphisme partiel entre
  $\varMA$ (la sous-structure de $\varMM$ engendrée par $\varPt$) et
  $\varMoM(\varMA)$.
\end{remark}

\begin{lemma}\label{lem.extend.part.elt.embed}
  Soient $\varMM,\varMN$ deux structures sur une signature $\Sign$, et
  $(\varPt,\varMoM)$ un plongement élémentaire partiel entre les deux
  structures. Alors il existe une extension $\varMN'$ de $\varMN$
  et un plongement élémentaire $\makeFunName{\varMoM'}{\varMM}{\varMN'}$ tel que
  $\varMoM'\restr{\varPt} = \varMoM$.
\end{lemma}

\begin{proof}
  On enrichit d'abord $\varMM$ et $\varMN$ en deux $\SignP{\varPt}$-structures,
  respectivement $\relev{\varMM}$ et $\relev{\varMN}$. Le premier enrichissement
  est canonique, et le second se fait par l'association
  $\Cc_{a} \mapsto\varMoM(a)$.

  Le fait que $(\varPt,\varMoM)$ est un plongement élémentaire partiel
  signifie en particulier que $\relev{\varMM}\eltEq\relev{\varMN}$ (en tant
  que $\SignP{\varPt}$-structures). D'après le théorème précédent, il existe
  donc une extension élémentaire commune aux deux structures.
  On trouve donc $\varMN'$ muni d'un plongement élémentaire
  $\makeFunName{\varMoM'}{\relev{\varMM}}{\varMN'}$ et d'un plongement
  élémentaire $\makeFunName{\varMoMB}{\relev{\varMN}}{\varMN'}$. Ainsi, quitte
  à remplacer $\varMoMB(\carMN)$ par $\carMN$, on obtient $\varMN'$ tel que
  $\varMN\subModElt\varMN'$ et un plongement élémentaire qui se restreint
  en $\varMoM$ sur $\varPt$.
\end{proof}

Le second outil que nous abordons pour étendre des structures et en construire
de nouvelles est celui des chaînes.

\begin{definition}[Chaîne de structures, chaîne élémentaire]
  Soit $\Sign$ une signature. Pour un ensemble totalement ordonné $(\indSet,<)$,
  on appelle chaîne de $\Sign$-structures une famille $(\varMM_i)_{i \in \indSet}$
  telle que
  \[\forall i,j \in \indSet, i < j \implies \varMM_i \subseteq\varMM_j\]
  On dit qu'une chaîne est élémentaire si la condition précédente est renforcée
  en
  \[\forall i,j \in \indSet, i < j \implies \varMM_i\subModElt\varMM_j\]
\end{definition}

L'intérêt d'une telle famille de structures est, assez naturellement, d'en
prendre l'union.

\begin{definition}[Union d'une chaîne de structures]
  Soit $\Sign$ une signature et $(\varMM_i)_{i \in \indSet}$ une chaîne
  de structure sur $(\indSet,<)$. On appelle union de cette chaîne la
  structure $\varMM$ définie par~:
  \begin{itemize}
  \item $\displaystyle\carMM \defeq \bigcup_{i \in \indSet}\varMM_i$
  \item pour tout symbole de fonction $\symbF$ d'arité $n$,
    $\displaystyle\toM{\symbF}\defeq
    \bigcup_{i \in \indSet}\toMP{\symbF}{\varMM_i}$
  \item pour tout symbole de relation $\symbR$ d'arité $n$,
    $\displaystyle\toM{\symbR}\defeq
    \bigcup_{i \in \indSet}\toMP{\symbR}{\varMM_i}$
  \end{itemize}
  on a alors naturellement, pour tout $i \in \indSet$,
  $\varMM_i\subseteq\varMM$. On note aussi $\bigcup_{i \in \indSet}\varMM_i$
  l'union de cette chaîne.
\end{definition}

\begin{remark}
  On peut souhaiter un point de vue plus structurel, dans lequel l'inclusion
  ne joue pas de rôle prépondérant. L'information essentielle est jouée par un
  plongement quelconque, et on pourrait donc souhaiter qu'une chaîne soit la
  donnée pour chaque $i < j$ d'un plongement élémentaire
  $\makeFunName{\varMoM_{i < j}}{\varMM_i}{\varMM_j}$
  vérifiant que $\varMoM_{j < k}\circ\varMoM_{i < j} = \varMoM_{i < k}$
  (ce que l'on pourrait considérer comme un certain foncteur de
  $(\indSet,\leq)$ vers la catégorie des structures munies des plongements).

  Dans un tel cadre, la construction est plus technique mais existe aussi~: on
  appelle une telle union une limite inductive, et on peut la construire en
  considérant des couples $(i,m)$ avec $m \in \carM{\varMM_i}$, puis en
  quotientant par la relation
  \[(i,m)\sim (j,m') \defeq
  \varMoM_{i < j}(m) = m' \lor \varMoM_{j < i}(m') = m\]
  Tout se passe encore de la même façon que précédemment~: cette limite
  inductive est naturellement munie d'une famille de plongements
  $\makeFunName{\varMoM_i}{\varMM_i}{\varMM}$ (qui est même universelle).
\end{remark}

\begin{remark}
  On peut caractériser assez simplement la sémantique dans l'union d'une
  chaîne~:
  \begin{itemize}
  \item $\displaystyle\bigcup_{i \in \indSet}\varMM_i,\vec a\satisf\varFF$ est
    vérifié exactement quand il existe $i\in\indSet$ tel que
    $\vec a \subseteq \varMM_i$ et $\varMM_i,\vec a \satisf\varFF$, ou de
    façon équivalence lorsque pour tout $i \in \indSet$ tel que
    $\vec a\subseteq\varMM_i$, $\varMM_i,\vec a \satisf\varFF$
  \item le calcul de $\toMP{\symbF}{\bigcup_{i \in \indSet}\varMM_i}(\vec a)$ se
    fait en considérant un (de façon équivalente, tous les) $i \in \indSet$
    tel que $\vec a \subseteq \varMM_i$, et est alors
    $\toMP{\symbF}{\varMM_i}(\vec a)$.
  \end{itemize}
  Le fait d'avoir des extensions de structures assure qu'on a bien une
  équivalence entre \og pour tout $i \in \indSet$ \fg et \og il existe
  $i \in \indSet$\fg dans les formulations données plus tôt.
\end{remark}

\begin{lemma}[Lemme de la chaîne élémentaire de Tarski \cite{Tarski1955}]
  Si $(\varMM_i)_{i \in \indSet}$ est une chaîne élémentaire, alors
  \[\forall i \in \indSet, \varMM_i\subModElt\bigcup_{i \in I}\varMM_i\]
\end{lemma}

\begin{proof}
  Notons $\varMM\defeq \bigcup_{i \in \indSet}\varMM_i$.
  On prouve par induction sur les formules $\varFF$ que pour tout
  $i \in \indSet$ et toute donnée de paramètres $\vec m$ dans $\varMM_i$,
  $\varMM_i,\vec m\satisf\varFF\iff\varMM,\vec m\satisf\varFF$~:
  \begin{itemize}
  \item Si $\varFF$ est atomique, alors le résultat est donné directement par
    le fait que $\varMM_i\subseteq\varMM$ (ou, pour une limite inductive, par
    le fait qu'on a un plongement).
  \item Si $\varFF$ est $\top$ ou $\bot$, le résultat est automatique.
  \item Si $\varFF = \lnot\varFP$, comme
    $\varMM_i,\vec m\satisf\varFP\iff\varMM,\vec m\satisf\varFP$, le résultat
    est encore vrai en passant à la négation.
  \item Les cas $\varFP\oplus\varFC$ où
    $\oplus\in\{\mathord{\lor},\mathord{\land},\mathord{\to}\}$ sont
    vrais en dépliant les définitions.
  \item Supposons que $\varFF = \exists x, \varFP$, et montrons le résultat
    par double implication~:
    \begin{itemize}
    \item[\fbox{$\Rightarrow$}] Supposons que
      $\varMM_i,\vec m\satisf\exists x, \varFP$. Alors on trouve
      $a\in\carM{\varMM_i}$ tel que $\varMM_i,(a,\vec m)\satisf\varFP$,
      donc par hypothèse d'induction $\varMM,(a,\vec m)\satisf\varFP$, donc
      $\varMM,\vec m\satisf\exists x, \varFP$.
    \item[\fbox{$\Leftarrow$}] Supposons que
      $\varMM,\vec m\satisf\exists x, \varFP$. Alors on trouve
      $j\in\indSet$ et $a \in \varMM_j$ tel que
      $\varMM,(a,\vec m)\satisf\varFP$. Comme $\indSet$ est totalement ordonné,
      on sait que $i$ et $j$ sont comparaibles. Quitte à prendre pour $j$ le
      plus grand des deux, on peut alors considérer que $\vec m \in\varMM_j$,
      et donc (par hypothèse d'induction) que
      $\varMM_j,(a,\vec m)\satisf\varFP$. On en déduit donc que
      $\varMM_j,\vec m\satisf\exists x, \varFP$, puis par inclusion élémentaire
      que $\varMM_i,\vec m\satisf\exists x, \varFP$.
    \end{itemize}
  \item Si $\varFF = \forall x, \varFP$, alors on peut se ramener aux cas
    précédents par les lois de De Morgan.
  \end{itemize}
  En particulier, si $\varFF$ est une proposition et $i\in\indSet$, on en
  déduit que
  \[\varMM_i\subModElt\varMM\qedhere\]
\end{proof}

En utilisant les chaînes élémentaires, on peut maintenant donner une version
plus forte du \cref{lem.extend.part.elt.embed} dans le cas d'un plongement
d'un modèle dans lui-même.

\begin{proposition}
  Soit $\Sign$ une signature, $\varMM$ une $\Sign$-structure et
  $(\varPt,\varMoM)$ un plongement élémentaire partiel de $\varMM$ dans
  elle-même. Alors il existe une extension élémentaire $\varMN$ de $\varMM$
  et $\makeFunNameIso{\relev{\varMoM}}{\varMN}{\varMN}$ étendant $\varMoM$.
\end{proposition}

\begin{proof}
  On construit une suite de structures à partir de $\varMM$.

  
  \begin{minipage}{0.3\textwidth}
    \begin{tikzpicture}
      \node (M0l) at (0,0) {$\varPt$};
      \node (M0r) at (2.3,0) {$\varMM_0$};
      \node (M1l) at (0,-2) {$\varMM_1$};
      \node (M1r) at (2.3,-2) {$\im(\varMoM_0)$};
      \node (M2l) at (0,-4) {$\im(\varMoMB_1)$};
      \node (M2r) at (2.3,-4) {$\varMM_2$};
      \node (M2lp) at (0,-4)[below] {$\vdots$};
      \node (M2rp) at (2.3,-4)[below] {$\vdots$};
      \path (M0l) --node[midway] {$\subseteq$} (M0r);
      \draw[->,>=latex]
      (M0l) --node[midway,left,dashed] {$\iota\circ\varMoM$} (M1l);
      \path (M1r) --node[midway] {$\supseteq$} (M1l);
      \draw[->,>=latex]
      (M1r) --node[midway,right,dashed] {$\iota\circ\varMoM_0^{-1}$} (M2r);
      \path (M2l) --node[midway] {$\subseteq$} (M2r);
      \draw[->,>=latex]
      (M1r) -- node[midway, right] {$\varMoM_0^{-1}$} (M0r);
      \draw[->,>=latex]
      (M2l) -- node[midway, left] {$\varMoMB_1^{-1}$} (M1l);
      \draw[->,>=latex,dashed]
      (M0r) -- node[midway,above left] {$\varMoM_0$} (M1l);
      \draw[->,>=latex,dashed]
      (M1l) -- node[midway,above right] {$\varMoMB_1$} (M2r);
    \end{tikzpicture}
  \end{minipage}
  \begin{minipage}{0.65\textwidth}
    On pose $\varMM_0 \defeq \varMM$.
    On sait d'après le \cref{lem.extend.part.elt.embed} qu'il existe une
    extension $\varMM_1$ et un plongement élémentaire
    $\makeFunName{\varMoM_0}{\varMM_0}{\varMM_1}$ qui se restreint en
    $\varMoM$ sur $\varPt$. On considère alors $\im(\varMoM_0)$, sur
    laquelle $\varMoM_0$ est bijective~: on a donc
    $\makeFunName{\varMoM_0^{-1}}{\im(\varMoM_0)}{\varMM_0}$, soit
    $\makeFunName{\iota\circ\varMoM_0^{-1}}{\im(\varMoM_0)}{\varMM_1}$ qui peut
    s'étendre par le \cref{lem.extend.part.elt.embed} en un plongement
    élémentaire $\makeFunName{\varMoMB_1}{\varMM_1}{\varMM_2}$ avec
    $\varMM_2$ une nouvelle structure. On remarque que
    $\varMoMB_1$ est l'inverse de $\varMoM_0$ quand on se restreint à
    l'image de $\varMoM_0$. On peut ainsi construire une suite de structures
    $\varMM_n$, une suite de plongements élémentaires $\varMoM_{2n}$ et une
    suite de plongements élémentaires $\varMoMB_{2n+1}$, de telle sorte que
    $\varMoM_{2n}\circ\varMoMB_{2n-1} = \id$ et
    $\varMoMB_{2n+1}\circ\varMoM_{2n}=\id$.
  \end{minipage}

  \vspace{5pt}
  Ainsi, en prenant l'union de cette chaîne sur $\bN$, on obtient une
  structure et deux morphismes
  \[\varMN\defeq\bigcup_{n \in \bN}\varMM_n\qquad
  \relev{\varMoM}\defeq \bigcup_{n \in \bN}\varMoM_{2n}\qquad
  \relev{\varMoMB}\defeq\bigcup_{n \in \bN}\varMoMB_{2n+1}\]
  donnant $\makeFunNameIso{\relev{\varMoM}}{\varMN}{\varMN}$ et
  $\relev{\varMoM}\restr{\varPt} = \varMoM$.
\end{proof}

\begin{corollary}
  Soit $\Sign$ une signature, et $\varMM$ une $\Sign$-structure. Pour tous
  $m,n \in \carMM$, si
  \[\forall \varFF(x)\in\Formula_1(\Sign), \varMM,m\satisf\varFF \iff
  \varMM,n\satisf\varFF\]
  alors il existe une extension élémentaire $\varMN$ de $\varMM$  et un
  automorphisme $\makeFunNameIso{\varMoM}{\varMN}{\varMN}$ tel que
  $\varMoM(m) = n$.
\end{corollary}

\begin{proof}
  $(\{m\},m\mapsto n)$ est un plongement élémentaire partiel de $\varMM$
  dans elle-même.
\end{proof}

\section{Parties définissables}
